% Generated by experiments/scripts/export_latex.py (export_latex_c12.py) -- do not edit by hand.
% Source: experiments/data/2-processed-tables-and-evaluations/revised-contract-12m/describe/describe.json
\begin{table}[htbp]
\centering
\caption[계약 그래프의 계정 종류.]{계약 그래프의 계정 종류예요. 종류 판별 기간에 다섯 토큰의 이벤트에는 0 주소가 아닌 서로 다른 주소가 25{,}537{,}280개 있어요. 그 가운데 17{,}477{,}779개는 트랜잭션(블록체인에 보내는 거래 요청)을 보냈고, 그 사실만으로 사용자 계정이에요. 이 표는 블록 513{,}216{,}506에서 \texttt{eth\_getCode}로 읽은 주소 6{,}230{,}270개의 종류를 판별해요. 여기에는 트랜잭션을 보낸 적이 없고 트랜잭션을 보낸 주소도 0 주소도 아닌 상대가 있는 모든 주소, 코호트의 구성원, 그리고 두 가지 감사 표본이 들어요. 이 집합 밖에서 트랜잭션을 보낸 적이 없는 주소는 상대가 모두 사용자 계정이어서, 계약 사이 화살표의 끝이 될 수 없어요. 스마트 계정은 사용자 계정이에요. Safe는 자기 문턱값과 주인들을 답해요. ERC-4337 계정은 표준 EntryPoint를 답하고, \texttt{validateUserOp} 셀렉터를 가지며 paymaster 셀렉터는 없어요. 프록시가 그 셀렉터를 가리는 계정 192{,}656개는 목록에 오른 코드 계열로 옮겼어요(코드 해시로 37{,}337개, 코드 해시와 구현으로 155{,}319개). 읽은 주소 가운데 524{,}454개는 ERC-1967 프록시, 503{,}516개는 ERC-1167 복제, 45{,}362개는 Safe 프록시예요. 트랜잭션을 보낸 주소까지 합치면, 주소 21{,}877{,}279개가 사용자 계정으로, 1{,}829{,}719개가 계약으로 판별돼요.}
\label{tab:contract-types}
\fitwidth{%
\begin{tabular}{llr}
\toprule
그 블록에서 읽은 코드 & 분류 & 주소 수 \\
\midrule
코드가 없는 외부 소유 계정 & 사용자 계정 & 2{,}909{,}509 \\
EIP-7702 위임이 있는 외부 소유 계정 & 사용자 계정 & 1{,}165{,}025 \\
Safe (문턱값과 주인들을 답함) & 사용자 계정 & 46{,}204 \\
ERC-4337 계정, 읽기 규칙으로 판별 & 사용자 계정 & 87{,}157 \\
ERC-4337 계정, 목록에 오른 코드 계열로 판별 & 사용자 계정 & 192{,}656 \\
그 밖의 모든 코드 & 계약 (그래프의 점) & 1{,}829{,}719 \\
\midrule
읽은 주소 & & 6{,}230{,}270 \\
\bottomrule
\end{tabular}
}
\end{table}
